%%%%%%%%%%%%%%%%%%%%%%%%%%%%%%%%%%%%%%%%%%%%%%%%%%%%%
%%   Atencao! A lista nao ordena automaticamente   %%
%%%%%%%%%%%%%%%%%%%%%%%%%%%%%%%%%%%%%%%%%%%%%%%%%%%%%
%
% Ordenada por tema, na sequência em que os símbolos aparecem no texto:
% métrica de fidelidade, desenho experimental, análise estatística, função de
% recompensa e, por fim, o símbolo grego. Contém apenas símbolos efetivamente
% empregados — as siglas das métricas (DF-F1, MF-F1, TCR, XSD-Val) figuram na
% Lista de Abreviaturas e Siglas, não aqui.

\begin{simbolos} \itemsep -1pt
	% Métrica de fidelidade topológica
	\item[$ R $] Multiconjunto \textit{direct-follows} do modelo de referência
	\item[$ C $] Multiconjunto \textit{direct-follows} do modelo candidato
	\item[$ P $] Precisão sobre o multiconjunto \textit{direct-follows}
	\item[$ R\!c $] Revocação sobre o multiconjunto \textit{direct-follows}
	\item[$ (a,b) $] Par de atividades em que $a$ é diretamente seguida por $b$

	% Desenho experimental
	\item[$ n $] Número de itens do conjunto de avaliação
	\item[$ k $] Número de repetições de geração por item

	% Análise estatística
	\item[$ p $] Valor-$p$ do teste de Wilcoxon pareado
	\item[$ p_{Holm} $] Valor-$p$ após correção de Holm para múltiplas comparações
	\item[$ r_{rb} $] Tamanho de efeito \textit{rank-biserial}

	% Função de recompensa da etapa opcional de aprendizado por reforço
	\item[$ r $] Recompensa composta
	\item[$ r' $] Recompensa restrita a componentes verificáveis sem referência
	\item[$ r_{\text{sint}} $] Componente de validade sintática da recompensa
	\item[$ r_{\text{sem}} $] Componente de adequação semântica da recompensa
	\item[$ r_{\text{topo}} $] Componente de coerência topológica da recompensa
	\item[$ r_{\text{comp}} $] Componente de economia de \textit{tokens} da recompensa

	% Símbolo grego
	\item[$ \alpha $] Algoritmo $\alpha$ de descoberta de processos
\end{simbolos}
